\section{F}

% fahren (to drive)
\begin{verbentry}{
  style = {acc},
  toc = {fahren (to drive)},
  name = {fahren},
  english = {to drive},
  type = {irregular, + Akkusativ},
  auxiliary = {haben/sein}
}
 
    
     \vspace{5pt}

    \hrule
    
    \vspace{5pt}

  \begin{tabular}{rl}
     \multicolumn{2}{l}{\textbf{PRÄSENS} }\\
    ich & \textbf{fahre}\\
    du & \textbf{f\"ahrst}\\
    er/sie/es & \textbf{f\"ahrt}\\
    wir & \textbf{fahren}\\
    ihr & \textbf{fahrt}\\
    sie/Sie & \textbf{fahren}\\
  \end{tabular}
  \hfill
     \begin{tabular}{rl}
    \multicolumn{2}{l}{\textbf{PRÄTERITUM} }\\
    ich & \textbf{fuhr}\\
    du & \textbf{fuhrst}\\
    er/sie/es & \textbf{fuhr}\\
    wir & \textbf{fuhren}\\
    ihr & \textbf{fuhrt}\\
    sie/Sie & \textbf{fuhren}\\
  \end{tabular}
  \hfill
  \begin{tabular}{rl}
     \multicolumn{2}{l}{\textbf{IMPERATIV} }\\
   & \textbf{fahr(e) (du)!}\\
   & \textbf{fahren wir!}\\
   & \textbf{fahrt (ihr)!}\\
   & \textbf{fahren Sie!}\\
   & \vspace{10pt}
  \end{tabular}

    \vspace{10pt}

 \begin{tabular}{rl}
     \multicolumn{2}{l}{\textbf{PERFEKT}}\\
  &  sein/haben + \textbf{gefahren}\\
  \end{tabular}
\hfill
   \begin{tabular}{rl}
   
    \multicolumn{2}{l}{\textbf{FUTURE I} }\\
  &  werden + \textbf{fahren}\\
 
  \end{tabular}
\hfill
   \begin{tabular}{rl}
      \multicolumn{2}{l}{\textbf{PLUSQUAMPERFEKT}}\\
   & sein/hatten + \textbf{gefahren}\\
  \end{tabular}
  
\vspace{10pt}

   \begin{tabular}{lll}
      \multicolumn{3}{l}{\textbf{MIT PRÄPOSITIONEN}}\\
& \textbf{fahren nach} + Ort & (to go to city / country)\\
&\textbf{fahren in} + Akk & (to go into)\\
&\textbf{fahren zu} + Dat & (to go to place / person)\\
&\textbf{fahren mit} + Dat & (to go with)\\
&\textbf{fahren über} + Akk & (to go via)\\
&\textbf{fahren von} + Akk  \textbf{nach} + Akk & (to go from ... to)\\
&\textbf{fahren durch} + Akk & (to go through)\\
\vspace{20pt}
  \end{tabular}
  \hfill
   \begin{tabular}{lll}
      \multicolumn{3}{l}{\textbf{TRENNBARE VERBEN}}\\
 & \textbf{ab$|$fahren} & (to depart)\\
 & \textbf{an$|$fahren} & (to start driving)\\
  & \textbf{aus$|$fahren} & (to drive out)\\
  & \textbf{ein$|$fahren} & (to enter / pull in)\\
  & \textbf{los$|$fahren} & (to set off)\\
     & \textbf{mit$|$fahren} & (to ride along)\\
     & \textbf{vor$|$fahren} & (to pull up)\\
     & \textbf{weg$|$fahren} & (to drive away)\\
      & \textbf{zurück$|$fahren} & (to bring back)\\
  \end{tabular}


  \vspace{-5pt}
\end{verbentry}

% fallen (to fall)
\begin{verbentry}{
  style = {default},
  toc = {fallen (to fall)},
  name = {fallen},
  english = {to fall},
  type = {irregular},
  auxiliary = {sein}
}
 
    
     \vspace{5pt}

    \hrule
    
    \vspace{5pt}

  \begin{tabular}{rl}
     \multicolumn{2}{l}{\textbf{PRÄSENS} }\\
    ich & \textbf{falle}\\
    du & \textbf{fällst}\\
    er/sie/es & \textbf{fällt}\\
    wir & \textbf{fallen}\\
    ihr & \textbf{fallt}\\
    sie/Sie & \textbf{fallen}\\
  \end{tabular}
  \hfill
     \begin{tabular}{rl}
    \multicolumn{2}{l}{\textbf{PRÄTERITUM} }\\
    ich & \textbf{fiel}\\
    du & \textbf{fielst}\\
    er/sie/es & \textbf{fiel}\\
    wir & \textbf{fielen}\\
    ihr & \textbf{fielt}\\
    sie/Sie & \textbf{fielen}\\
  \end{tabular}
  \hfill
  \begin{tabular}{rl}
     \multicolumn{2}{l}{\textbf{IMPERATIV} }\\
   & \textbf{fall(e) (du)!}\\
   & \textbf{fallen wir!}\\
   & \textbf{fallt (ihr)!}\\
   & \textbf{fallen Sie!}\\
   & \vspace{10pt}
  \end{tabular}

    \vspace{10pt}

 \begin{tabular}{rl}
     \multicolumn{2}{l}{\textbf{PERFEKT}}\\
    & sein + \textbf{gefallen}\\
  \end{tabular}
\hfill
   \begin{tabular}{rl}
   
    \multicolumn{2}{l}{\textbf{FUTURE I} }\\
    &werden + \textbf{fallen}\\
 
  \end{tabular}
\hfill
   \begin{tabular}{rl}
      \multicolumn{2}{l}{\textbf{PLUSQUAMPERFEKT}}\\
    &waren + \textbf{gefallen}\\
  \end{tabular}

   \vspace{10pt}

   \begin{tabular}{lll}
      \multicolumn{3}{l}{\textbf{MIT PRÄPOSITIONEN}}\\
& \textbf{fallen in} + Akk & (to fall into)\\
& \textbf{fallen auf} + Akk & (to fall onto)\\
&\textbf{fallen von} + Dat & (to fall from)\\
&\textbf{fallen über} + Akk & (to fall over)\\
&\textbf{fallen an} + Akk & (to accrue / arise costs or interest)\\
\vspace{10pt} 
\end{tabular}
  \hfill
   \begin{tabular}{lll}
      \multicolumn{3}{l}{\textbf{TRENNBARE VERBEN}}\\
 & \textbf{ab$|$fallen} & (to fall / drop off)\\
  & \textbf{auf$|$fallen} & (to stand out)\\
   & \textbf{aus$|$fallen} & (to be cancelled / to fail)\\
   & \textbf{ein$|$fallen} & (to occur to someone)\\
    & \textbf{um$|$fallen} & (to fall over)\\
       & \textbf{zurück$|$fallen} & (to fall behind)\\
  \end{tabular}


  \vspace{-5pt}
\end{verbentry}

% fangen (to catch / capture)
\begin{verbentry}{
  style = {acc},
  toc = {fangen (to catch / capture)},
  name = {fangen},
  english = {to catch / capture},
  type = {irregular, + Akkusativ},
  auxiliary = {haben}
}
 
    
     \vspace{5pt}

    \hrule
    
    \vspace{5pt}

  \begin{tabular}{rl}
     \multicolumn{2}{l}{\textbf{PRÄSENS} }\\
    ich & \textbf{fange}\\
    du & \textbf{fängst}\\
    er/sie/es & \textbf{fängt}\\
    wir & \textbf{fangen}\\
    ihr & \textbf{fangt}\\
    sie/Sie & \textbf{fangen}\\
  \end{tabular}
  \hfill
     \begin{tabular}{rl}
    \multicolumn{2}{l}{\textbf{PRÄTERITUM} }\\
    ich & \textbf{fing}\\
    du & \textbf{fingst}\\
    er/sie/es & \textbf{fing}\\
    wir & \textbf{fingen}\\
    ihr & \textbf{fingt}\\
    sie/Sie & \textbf{fingen}\\
  \end{tabular}
  \hfill
  \begin{tabular}{rl}
     \multicolumn{2}{l}{\textbf{IMPERATIV} }\\
   & \textbf{fang(e) (du)!}\\
   & \textbf{fangen wir!}\\
   & \textbf{fangt (ihr)!}\\
   & \textbf{fangen Sie!}\\
   & \vspace{10pt}
  \end{tabular}

    \vspace{10pt}

 \begin{tabular}{rl}
     \multicolumn{2}{l}{\textbf{PERFEKT}}\\
  &  haben + \textbf{gefangen}\\
  \end{tabular}
\hfill
   \begin{tabular}{rl}
   
    \multicolumn{2}{l}{\textbf{FUTURE I} }\\
  &  werden + \textbf{fangen}\\
 
  \end{tabular}
\hfill
   \begin{tabular}{rl}
      \multicolumn{2}{l}{\textbf{PLUSQUAMPERFEKT}}\\
  &  hatten + \textbf{gefangen}\\
  \end{tabular}
  
  \vspace{10pt}

   \begin{tabular}{lll}
      \multicolumn{3}{l}{\textbf{MIT PRÄPOSITIONEN}}\\
      & \textbf{fangen nach} + Dat & (to reach for / try to catch)\\
&\textbf{anfangen mit} + Dat & (to begin with)\\
&\textbf{anfangen bei} + Dat & (to begin at (point / person))\\
& \textbf{einfangen in} + Akk & (to capture in)\\
& \textbf{abfangen an} + Dat & (to intercept at)
  \end{tabular}
  \hfill
   \begin{tabular}{lll}
      \multicolumn{3}{l}{\textbf{TRENNBARE VERBEN}}\\
 & \textbf{an$|$fangen} & (to begin)\\
 & \textbf{auf$|$fangen} & (to catch / absorb / cushion)\\
 & \textbf{ein$|$fangen} & (to catch / capture)\\
  & \textbf{ab$|$fangen} & (to intercept)\\
  & \textbf{weg$|$fangen} & (to snatch away)\\
  & \textbf{über$|$fangen} & (to cover / overtake)
  \end{tabular}


  \vspace{-5pt}
\end{verbentry}

% faszinieren (to fascinate)
\begin{verbentry}{
  style = {acc},
  toc = {faszinieren (to fascinate)},
  name = {faszinieren},
  english = {to fascinate},
  type = {regular, + Akkusativ},
  auxiliary = {haben}
}


\vspace{5pt}
\hrule
\vspace{5pt}

\begin{tabular}{rl}
\multicolumn{2}{l}{\textbf{PRÄSENS}}\\
ich & \textbf{fasziniere}\\
du & \textbf{faszinierst}\\
er/sie/es & \textbf{fasziniert}\\
wir & \textbf{faszinieren}\\
ihr & \textbf{fasziniert}\\
sie/Sie & \textbf{faszinieren}\\
\end{tabular}
\hfill
\begin{tabular}{rl}
\multicolumn{2}{l}{\textbf{PRÄTERITUM}}\\
ich & \textbf{faszinierte}\\
du & \textbf{fasziniertest}\\
er/sie/es & \textbf{faszinierte}\\
wir & \textbf{faszinierten}\\
ihr & \textbf{fasziniertet}\\
sie/Sie & \textbf{faszinierten}\\
\end{tabular}
\hfill
\begin{tabular}{rl}
\multicolumn{2}{l}{\textbf{IMPERATIV}}\\
& \textbf{fasziniere (du)!}\\
& \textbf{faszinieren wir!}\\
& \textbf{fasziniert (ihr)!}\\
& \textbf{faszinieren Sie!}\\
\end{tabular}

\vspace{10pt}

\begin{tabular}{rl}
\multicolumn{2}{l}{\textbf{PERFEKT}}\\
& haben + \textbf{fasziniert}\\
\end{tabular}
\hfill
\begin{tabular}{rl}
\multicolumn{2}{l}{\textbf{FUTURE I}}\\
& werden + \textbf{faszinieren}\\
\end{tabular}
\hfill
\begin{tabular}{rl}
\multicolumn{2}{l}{\textbf{PLUSQUAMPERFEKT}}\\
& hatten + \textbf{fasziniert}\\
\end{tabular}

\vspace{10pt}

\begin{tabular}{lll}
\multicolumn{3}{l}{\textbf{MIT PRÄPOSITIONEN}}\\
& \textbf{---} & \\
\end{tabular}
\hfill
\begin{tabular}{lll}
\multicolumn{3}{l}{\textbf{TRENNBARE VERBEN}}\\
& \textbf{---} & \\
\end{tabular}

\vspace{-5pt}
\end{verbentry}

% fehlen (to be missing)
\begin{verbentry}{
  style = {default},
  toc = {fehlen (to be missing)},
  name = {fehlen},
  english = {to be missing},
  type = {regular},
  auxiliary = {haben}
}
 
    
     \vspace{5pt}

    \hrule
    
    \vspace{5pt}

  \begin{tabular}{rl}
     \multicolumn{2}{l}{\textbf{PRÄSENS} }\\
    ich & \textbf{fehle}\\
    du & \textbf{fehlst}\\
    er/sie/es & \textbf{fehlt}\\
    wir & \textbf{fehlen}\\
    ihr & \textbf{fehlt}\\
    sie/Sie & \textbf{fehlen}\\
  \end{tabular}
  \hfill
     \begin{tabular}{rl}
    \multicolumn{2}{l}{\textbf{PRÄTERITUM} }\\
    ich & \textbf{fehlte}\\
    du & \textbf{fehltest}\\
    er/sie/es & \textbf{fehlte}\\
    wir & \textbf{fehlten}\\
    ihr & \textbf{fehltet}\\
    sie/Sie & \textbf{fehlten}\\
  \end{tabular}
  \hfill
  \begin{tabular}{rl}
     \multicolumn{2}{l}{\textbf{IMPERATIV} }\\
   & \textbf{fehl(e) (du)!}\\
   & \textbf{fehlen wir!}\\
   & \textbf{fehlt (ihr)!}\\
   & \textbf{fehlen Sie!}\\
   & \vspace{10pt}
  \end{tabular}

    \vspace{10pt}

 \begin{tabular}{rl}
     \multicolumn{2}{l}{\textbf{PERFEKT}}\\
    &haben + \textbf{gefehlt}\\
  \end{tabular}
\hfill
   \begin{tabular}{rl}
   
    \multicolumn{2}{l}{\textbf{FUTURE I} }\\
   & werden + \textbf{fehlen}\\
 
  \end{tabular}
\hfill
   \begin{tabular}{rl}
      \multicolumn{2}{l}{\textbf{PLUSQUAMPERFEKT}}\\
    &hatten + \textbf{gefehlt}\\
  \end{tabular}
  
  \vspace{10pt}

\begin{tabular}{lll}
\multicolumn{3}{l}{\textbf{MIT PRÄPOSITIONEN}}\\
& \textbf{---} & \\
\end{tabular}
\hfill
\begin{tabular}{lll}
\multicolumn{3}{l}{\textbf{TRENNBARE VERBEN}}\\
& \textbf{---} & \\
\end{tabular}

  \vspace{-5pt}
\end{verbentry}

% feiern (to celebrate)
\begin{verbentry}{
  style = {acc},
  toc = {feiern (to celebrate)},
  name = {feiern},
  english = {to celebrate},
  type = {regular, + Akkusativ},
  auxiliary = {haben}
}
 
    
     \vspace{5pt}

    \hrule
    
    \vspace{5pt}

  \begin{tabular}{rl}
     \multicolumn{2}{l}{\textbf{PRÄSENS} }\\
    ich & \textbf{feiere}\\
    du & \textbf{feierst}\\
    er/sie/es & \textbf{feiert}\\
    wir & \textbf{feiern}\\
    ihr & \textbf{feiert}\\
    sie/Sie & \textbf{feiern}\\
  \end{tabular}
  \hfill
     \begin{tabular}{rl}
    \multicolumn{2}{l}{\textbf{PRÄTERITUM} }\\
    ich & \textbf{feierte}\\
    du & \textbf{feiertest}\\
    er/sie/es & \textbf{feierte}\\
    wir & \textbf{feierten}\\
    ihr & \textbf{feiertet}\\
    sie/Sie & \textbf{feierten}\\
  \end{tabular}
  \hfill
  \begin{tabular}{rl}
     \multicolumn{2}{l}{\textbf{IMPERATIV} }\\
   & \textbf{feier(e) (du)!}\\
   & \textbf{feiern wir!}\\
   & \textbf{feiert (ihr)!}\\
   & \textbf{feiern Sie!}\\
   & \vspace{10pt}
  \end{tabular}

    \vspace{10pt}

 \begin{tabular}{rl}
     \multicolumn{2}{l}{\textbf{PERFEKT}}\\
 &   haben + \textbf{gefeiert}\\
  \end{tabular}
\hfill
   \begin{tabular}{rl}
   
    \multicolumn{2}{l}{\textbf{FUTURE I} }\\
  &  werden + \textbf{feiern}\\
 
  \end{tabular}
\hfill
   \begin{tabular}{rl}
      \multicolumn{2}{l}{\textbf{PLUSQUAMPERFEKT}}\\
  &  hatten + \textbf{gefeiert}\\
  \end{tabular}
  
  
\vspace{10pt}

\begin{tabular}{lll}
\multicolumn{3}{l}{\textbf{MIT PRÄPOSITIONEN}}\\
& \textbf{---} & \\
\end{tabular}
\hfill
\begin{tabular}{lll}
\multicolumn{3}{l}{\textbf{TRENNBARE VERBEN}}\\
& \textbf{---} & \\
\end{tabular}
\vspace{-5pt}
\end{verbentry}

% finanzieren (to finance / fund)
\begin{verbentry}{
  style = {default},
  toc = {finanzieren (to finance / fund)},
  name = {finanzieren},
  english = {to finance / fund},
  type = {regular},
  auxiliary = {haben}
}


\vspace{5pt}
\hrule
\vspace{5pt}

\begin{tabular}{rl}
\multicolumn{2}{l}{\textbf{PRÄSENS}}\\
ich & \textbf{finanziere}\\
du & \textbf{finanzierst}\\
er/sie/es & \textbf{finanziert}\\
wir & \textbf{finanzieren}\\
ihr & \textbf{finanziert}\\
sie/Sie & \textbf{finanzieren}\\
\end{tabular}
\hfill
\begin{tabular}{rl}
\multicolumn{2}{l}{\textbf{PRÄTERITUM}}\\
ich & \textbf{finanzierte}\\
du & \textbf{finanziertest}\\
er/sie/es & \textbf{finanzierte}\\
wir & \textbf{finanzierten}\\
ihr & \textbf{finanziertet}\\
sie/Sie & \textbf{finanzierten}\\
\end{tabular}
\hfill
\begin{tabular}{rl}
\multicolumn{2}{l}{\textbf{IMPERATIV}}\\
& \textbf{finanziere (du)!}\\
& \textbf{finanzieren wir!}\\
& \textbf{finanziert (ihr)!}\\
& \textbf{finanzieren Sie!}\\
\end{tabular}

\vspace{10pt}

\begin{tabular}{rl}
\multicolumn{2}{l}{\textbf{PERFEKT}}\\
& haben + \textbf{finanziert}\\
\end{tabular}
\hfill
\begin{tabular}{rl}
\multicolumn{2}{l}{\textbf{FUTURE I}}\\
& werden + \textbf{finanzieren}\\
\end{tabular}
\hfill
\begin{tabular}{rl}
\multicolumn{2}{l}{\textbf{PLUSQUAMPERFEKT}}\\
& hatten + \textbf{finanziert}\\
\end{tabular}

\vspace{10pt}

\begin{tabular}{lll}
\multicolumn{3}{l}{\textbf{MIT PRÄPOSITIONEN}}\\
& \textbf{finanzieren mit} + Dat & (to finance with)\\
& \textbf{finanzieren durch} + Akk & (to finance through / by)\\
\end{tabular}
\hfill
\begin{tabular}{lll}
\multicolumn{3}{l}{\textbf{TRENNBARE VERBEN}}\\
& \textbf{---} & \\
\end{tabular}
\vspace{-5pt}
\end{verbentry}

% finden (to find)
\begin{verbentry}{
  style = {acc},
  toc = {finden (to find)},
  name = {finden},
  english = {to find},
  type = {irregular, + Akkusativ},
  auxiliary = {haben}
}
 
    
     \vspace{5pt}

    \hrule
    
    \vspace{5pt}

  \begin{tabular}{rl}
     \multicolumn{2}{l}{\textbf{PRÄSENS} }\\
    ich & \textbf{finde}\\
    du & \textbf{findest}\\
    er/sie/es & \textbf{findet}\\
    wir & \textbf{finden}\\
    ihr & \textbf{findet}\\
    sie/Sie & \textbf{finden}\\
  \end{tabular}
  \hfill
     \begin{tabular}{rl}
    \multicolumn{2}{l}{\textbf{PRÄTERITUM} }\\
    ich & \textbf{fand}\\
    du & \textbf{fandest}\\
    er/sie/es & \textbf{fand}\\
    wir & \textbf{fanden}\\
    ihr & \textbf{fandet}\\
    sie/Sie & \textbf{fanden}\\
  \end{tabular}
  \hfill
  \begin{tabular}{rl}
     \multicolumn{2}{l}{\textbf{IMPERATIV} }\\
   & \textbf{find(e) (du)!}\\
   & \textbf{finden wir!}\\
   & \textbf{findet (ihr)!}\\
   & \textbf{finden Sie!}\\
   & \vspace{10pt}
  \end{tabular}

    \vspace{10pt}

 \begin{tabular}{rl}
     \multicolumn{2}{l}{\textbf{PERFEKT}}\\
  &  haben + \textbf{gefunden}\\
  \end{tabular}
\hfill
   \begin{tabular}{rl}
   
    \multicolumn{2}{l}{\textbf{FUTURE I} }\\
  &  werden + \textbf{finden}\\
 
  \end{tabular}
\hfill
   \begin{tabular}{rl}
      \multicolumn{2}{l}{\textbf{PLUSQUAMPERFEKT}}\\
   & hatten + \textbf{gefunden}\\
  \end{tabular}
  
   \vspace{10pt}

\begin{tabular}{lll}
      \multicolumn{3}{l}{\textbf{MIT PRÄPOSITIONEN}}\\
& \textbf{finden für} + Akk & (to find suitable for)\\
& \textbf{finden zu} + Dat & (to find to / tend to)\\
& \textbf{finden in} + Dat & (to find in / inside)\\
\end{tabular}
  \hfill
\begin{tabular}{lll}
      \multicolumn{3}{l}{\textbf{TRENNBARE VERBEN}}\\
 & \textbf{wieder$|$finden} & (to recover / find again)\\
 & \textbf{statt$|$finden} & (to take place / happen)\\
\vspace{10pt}
  \end{tabular}

  \vspace{-5pt}
\end{verbentry}

% fliegen (to fly)
\begin{verbentry}{
  style = {default},
  toc = {fliegen (to fly)},
  name = {fliegen},
  english = {to fly},
  type = {irregular},
  auxiliary = {sein}
}
 
    
     \vspace{5pt}

    \hrule
    
    \vspace{5pt}

  \begin{tabular}{rl}
     \multicolumn{2}{l}{\textbf{PRÄSENS} }\\
    ich & \textbf{fliege}\\
    du & \textbf{fliegst}\\
    er/sie/es & \textbf{fliegt}\\
    wir & \textbf{fliegen}\\
    ihr & \textbf{fliegt}\\
    sie/Sie & \textbf{fliegen}\\
  \end{tabular}
  \hfill
     \begin{tabular}{rl}
    \multicolumn{2}{l}{\textbf{PRÄTERITUM} }\\
    ich & \textbf{flog}\\
    du & \textbf{flogst}\\
    er/sie/es & \textbf{flog}\\
    wir & \textbf{flogen}\\
    ihr & \textbf{flogt}\\
    sie/Sie & \textbf{flogen}\\
  \end{tabular}
  \hfill
  \begin{tabular}{rl}
     \multicolumn{2}{l}{\textbf{IMPERATIV} }\\
   & \textbf{flieg(e) (du)!}\\
   & \textbf{fliegen wir!}\\
   & \textbf{fliegt (ihr)!}\\
   & \textbf{fliegen Sie!}\\
   & \vspace{10pt}
  \end{tabular}

    \vspace{10pt}

 \begin{tabular}{rl}
     \multicolumn{2}{l}{\textbf{PERFEKT}}\\
   & sein + \textbf{geflogen}\\
  \end{tabular}
\hfill
   \begin{tabular}{rl}
    \multicolumn{2}{l}{\textbf{FUTURE I} }\\
   & werden + \textbf{fliegen}\\
  \end{tabular}
\hfill
   \begin{tabular}{rl}
      \multicolumn{2}{l}{\textbf{PLUSQUAMPERFEKT}}\\
   & waren + \textbf{geflogen}\\
  \end{tabular}

   \vspace{10pt}

   \begin{tabular}{lll}
      \multicolumn{3}{l}{\textbf{MIT PRÄPOSITIONEN}}\\
& \textbf{fliegen nach} + Dat & (to fly to – cities/countries)\\
& \textbf{fliegen in} + Akk & (to fly into)\\
& \textbf{fliegen über} + Akk & (to fly over)\\
& \textbf{fliegen von} + Dat & (to fly from)\\
   \end{tabular}
  \hfill
   \begin{tabular}{lll}
      \multicolumn{3}{l}{\textbf{TRENNBARE VERBEN}}\\
 & \textbf{ab$|$fliegen} & (to depart by plane)\\
 & \textbf{an$|$fliegen} & (to approach by flying)\\
 & \textbf{ein$|$fliegen} & (to fly in)\\
 & \textbf{weg$|$fliegen} & (to fly away)\\
 & \textbf{über$|$fliegen} & (to fly over)\\
  \end{tabular}

\vspace{-5pt}
\end{verbentry}

% fliehen (to flee)
\begin{verbentry}{
  style = {default},
  toc = {fliehen (to flee)},
  name = {fliehen},
  english = {to flee},
  type = {irregular},
  auxiliary = {sein}
}
 
    
     \vspace{5pt}

    \hrule
    
    \vspace{5pt}

  \begin{tabular}{rl}
     \multicolumn{2}{l}{\textbf{PRÄSENS} }\\
    ich & \textbf{fliehe}\\
    du & \textbf{fliehst}\\
    er/sie/es & \textbf{flieht}\\
    wir & \textbf{fliehen}\\
    ihr & \textbf{flieht}\\
    sie/Sie & \textbf{fliehen}\\
  \end{tabular}
  \hfill
     \begin{tabular}{rl}
    \multicolumn{2}{l}{\textbf{PRÄTERITUM} }\\
    ich & \textbf{floh}\\
    du & \textbf{flohst}\\
    er/sie/es & \textbf{floh}\\
    wir & \textbf{flohen}\\
    ihr & \textbf{floht}\\
    sie/Sie & \textbf{flohen}\\
  \end{tabular}
  \hfill
  \begin{tabular}{rl}
     \multicolumn{2}{l}{\textbf{IMPERATIV} }\\
   & \textbf{flieh(e) (du)!}\\
   & \textbf{fliehen wir!}\\
   & \textbf{flieht (ihr)!}\\
   & \textbf{fliehen Sie!}\\
   & \vspace{10pt}
  \end{tabular}

    \vspace{10pt}

 \begin{tabular}{rl}
     \multicolumn{2}{l}{\textbf{PERFEKT}}\\
   & sein + \textbf{geflohen}\\
  \end{tabular}
\hfill
   \begin{tabular}{rl}
    \multicolumn{2}{l}{\textbf{FUTURE I} }\\
   & werden + \textbf{fliehen}\\
  \end{tabular}
\hfill
   \begin{tabular}{rl}
      \multicolumn{2}{l}{\textbf{PLUSQUAMPERFEKT}}\\
   & waren + \textbf{geflohen}\\
  \end{tabular}

   \vspace{10pt}

   \begin{tabular}{lll}
      \multicolumn{3}{l}{\textbf{MIT PRÄPOSITIONEN}}\\
& \textbf{fliehen vor} + Dat & (to flee from)\\
& \textbf{fliehen aus} + Dat & (to flee out of / from)\\
\end{tabular}
  \hfill
   \begin{tabular}{lll}
      \multicolumn{3}{l}{\textbf{TRENNBARE VERBEN}}\\
 & \textbf{ent$|$fliehen} & (to escape from)\\
\vspace{10pt}
  \end{tabular}

  \vspace{-5pt}
\end{verbentry}

% folgen (to follow)
\begin{verbentry}{
  style = {dat},
  toc = {folgen (to follow)},
  name = {folgen},
  english = {to follow},
  type = {regular},
  auxiliary = {sein}
}
 
    
     \vspace{5pt}

    \hrule
    
    \vspace{5pt}

  \begin{tabular}{rl}
     \multicolumn{2}{l}{\textbf{PRÄSENS} }\\
    ich & \textbf{folge}\\
    du & \textbf{folgst}\\
    er/sie/es & \textbf{folgt}\\
    wir & \textbf{folgen}\\
    ihr & \textbf{folgt}\\
    sie/Sie & \textbf{folgen}\\
  \end{tabular}
  \hfill
     \begin{tabular}{rl}
    \multicolumn{2}{l}{\textbf{PRÄTERITUM} }\\
    ich & \textbf{folgte}\\
    du & \textbf{folgtest}\\
    er/sie/es & \textbf{folgte}\\
    wir & \textbf{folgten}\\
    ihr & \textbf{folgtet}\\
    sie/Sie & \textbf{folgten}\\
  \end{tabular}
  \hfill
  \begin{tabular}{rl}
     \multicolumn{2}{l}{\textbf{IMPERATIV} }\\
   & \textbf{folg(e) (du)!}\\
   & \textbf{folgen wir!}\\
   & \textbf{folgt (ihr)!}\\
   & \textbf{folgen Sie!}\\
   & \vspace{10pt}
  \end{tabular}

    \vspace{10pt}

 \begin{tabular}{rl}
     \multicolumn{2}{l}{\textbf{PERFEKT}}\\
   & sein + \textbf{gefolgt}\\
  \end{tabular}
\hfill
   \begin{tabular}{rl}
   
    \multicolumn{2}{l}{\textbf{FUTURE I} }\\
   & werden + \textbf{folgen}\\
 
  \end{tabular}
\hfill
   \begin{tabular}{rl}
      \multicolumn{2}{l}{\textbf{PLUSQUAMPERFEKT}}\\
   & waren + \textbf{gefolgt}\\
  \end{tabular}

   \vspace{10pt}

   \begin{tabular}{lll}
      \multicolumn{3}{l}{\textbf{MIT PRÄPOSITIONEN}}\\
& \textbf{folgen auf} + Akk & (to follow after)\\
\vspace{10pt}
\end{tabular}
  \hfill
   \begin{tabular}{lll}
      \multicolumn{3}{l}{\textbf{TRENNBARE VERBEN}}\\
 & \textbf{nach$|$folgen} & (to succeed / follow after)\\
 & \textbf{ver$|$folgen} & (not separable: to pursue)\\
  \end{tabular}


  \vspace{-5pt}
\end{verbentry}

% fragen (to ask)
\begin{verbentry}{
  style = {acc},
  toc = {fragen (to ask)},
  name = {fragen},
  english = {to ask},
  type = {regular, + Akkusativ},
  auxiliary = {haben}
}
 
    
     \vspace{5pt}

    \hrule
    
    \vspace{5pt}

  \begin{tabular}{rl}
     \multicolumn{2}{l}{\textbf{PRÄSENS} }\\
    ich & \textbf{frage}\\
    du & \textbf{fragst}\\
    er/sie/es & \textbf{fragt}\\
    wir & \textbf{fragen}\\
    ihr & \textbf{fragt}\\
    sie/Sie & \textbf{fragen}\\
  \end{tabular}
  \hfill
     \begin{tabular}{rl}
    \multicolumn{2}{l}{\textbf{PRÄTERITUM} }\\
    ich & \textbf{fragte}\\
    du & \textbf{fragtest}\\
    er/sie/es & \textbf{fragte}\\
    wir & \textbf{fragten}\\
    ihr & \textbf{fragtet}\\
    sie/Sie & \textbf{fragten}\\
  \end{tabular}
  \hfill
  \begin{tabular}{rl}
     \multicolumn{2}{l}{\textbf{IMPERATIV} }\\
   & \textbf{frag(e) (du)!}\\
   & \textbf{fragen wir!}\\
   & \textbf{fragt (ihr)!}\\
   & \textbf{fragen Sie!}\\
   & \vspace{10pt}
  \end{tabular}

    \vspace{10pt}

 \begin{tabular}{rl}
     \multicolumn{2}{l}{\textbf{PERFEKT}}\\
   & haben + \textbf{gefragt}\\
  \end{tabular}
\hfill
   \begin{tabular}{rl}
   
    \multicolumn{2}{l}{\textbf{FUTURE I} }\\
   & werden + \textbf{fragen}\\
 
  \end{tabular}
\hfill
   \begin{tabular}{rl}
      \multicolumn{2}{l}{\textbf{PLUSQUAMPERFEKT}}\\
   & hatten + \textbf{gefragt}\\
  \end{tabular}
  
  
\vspace{10pt}

\begin{tabular}{lll}
\multicolumn{3}{l}{\textbf{MIT PRÄPOSITIONEN}}\\
& \textbf{---} & \\
\end{tabular}
\hfill
\begin{tabular}{lll}
\multicolumn{3}{l}{\textbf{TRENNBARE VERBEN}}\\
& \textbf{---} & \\
\end{tabular}
\vspace{-5pt}
\end{verbentry}

% fressen (to eat — animals)
\begin{verbentry}{
  style = {default},
  toc = {fressen (to eat — animals)},
  name = {fressen},
  english = {to eat — animals},
  type = {irregular},
  auxiliary = {haben}
}
 
    
     \vspace{5pt}

    \hrule
    
    \vspace{5pt}

  \begin{tabular}{rl}
     \multicolumn{2}{l}{\textbf{PRÄSENS} }\\
    ich & \textbf{fresse}\\
    du & \textbf{frisst}\\
    er/sie/es & \textbf{frisst}\\
    wir & \textbf{fressen}\\
    ihr & \textbf{fresst}\\
    sie/Sie & \textbf{fressen}\\
  \end{tabular}
  \hfill
     \begin{tabular}{rl}
    \multicolumn{2}{l}{\textbf{PRÄTERITUM} }\\
    ich & \textbf{fraß}\\
    du & \textbf{fraßest}\\
    er/sie/es & \textbf{fraß}\\
    wir & \textbf{fraßen}\\
    ihr & \textbf{fraßt}\\
    sie/Sie & \textbf{fraßen}\\
  \end{tabular}
  \hfill
  \begin{tabular}{rl}
     \multicolumn{2}{l}{\textbf{IMPERATIV} }\\
   & \textbf{friss (du)!}\\
   & \textbf{fressen wir!}\\
   & \textbf{fresst (ihr)!}\\
   & \textbf{fressen Sie!}\\
   & \vspace{10pt}
  \end{tabular}

    \vspace{10pt}

 \begin{tabular}{rl}
     \multicolumn{2}{l}{\textbf{PERFEKT}}\\
   & haben + \textbf{gefressen}\\
  \end{tabular}
\hfill
   \begin{tabular}{rl}
    \multicolumn{2}{l}{\textbf{FUTURE I} }\\
   & werde + \textbf{fressen}\\
  \end{tabular}
\hfill
   \begin{tabular}{rl}
      \multicolumn{2}{l}{\textbf{PLUSQUAMPERFEKT}}\\
   & hatten + \textbf{gefressen}\\
  \end{tabular}

   \vspace{10pt}

   \begin{tabular}{lll}
      \multicolumn{3}{l}{\textbf{MIT PRÄPOSITIONEN}}\\
& \textbf{fressen von} + Dat & (to eat from / of)\\
& \textbf{fressen aus} + Dat & (to eat from)\\
\end{tabular}
  \hfill
   \begin{tabular}{lll}
      \multicolumn{3}{l}{\textbf{TRENNBARE VERBEN}}\\
 & \textbf{auf$|$fressen} & (to eat up)\\
 & \textbf{an$|$fressen} & (to nibble at / eat into)\\
\vspace{10pt}
  \end{tabular}

  \vspace{-5pt}
\end{verbentry}

% fügen (to add / insert)
\begin{verbentry}{
  style = {default},
  toc = {fügen (to add / insert)},
  name = {fügen},
  english = {to add / insert},
  type = {regular},
  auxiliary = {haben}
}


\vspace{5pt}
\hrule
\vspace{5pt}

\begin{tabular}{rl}
\multicolumn{2}{l}{\textbf{PRÄSENS}}\\
ich & \textbf{füge}\\
du & \textbf{fügst}\\
er/sie/es & \textbf{fügt}\\
wir & \textbf{fügen}\\
ihr & \textbf{fügt}\\
sie/Sie & \textbf{fügen}\\
\end{tabular}
\hfill
\begin{tabular}{rl}
\multicolumn{2}{l}{\textbf{PRÄTERITUM}}\\
ich & \textbf{fügte}\\
du & \textbf{fügtest}\\
er/sie/es & \textbf{fügte}\\
wir & \textbf{fügten}\\
ihr & \textbf{fügtet}\\
sie/Sie & \textbf{fügten}\\
\end{tabular}
\hfill
\begin{tabular}{rl}
\multicolumn{2}{l}{\textbf{IMPERATIV}}\\
& \textbf{füg(e) (du)!}\\
& \textbf{fügen wir!}\\
& \textbf{fügt (ihr)!}\\
& \textbf{fügen Sie!}\\
\end{tabular}

\vspace{10pt}

\begin{tabular}{rl}
\multicolumn{2}{l}{\textbf{PERFEKT}}\\
& haben + \textbf{gefügt}\\
\end{tabular}
\hfill
\begin{tabular}{rl}
\multicolumn{2}{l}{\textbf{FUTURE I}}\\
& werden + \textbf{fügen}\\
\end{tabular}
\hfill
\begin{tabular}{rl}
\multicolumn{2}{l}{\textbf{PLUSQUAMPERFEKT}}\\
& hatten + \textbf{gefügt}\\
\end{tabular}

\vspace{10pt}

\begin{tabular}{lll}
\multicolumn{3}{l}{\textbf{MIT PRÄPOSITIONEN}}\\
& \textbf{sich fügen in} + Akk & (to submit to)\\
\end{tabular}
\hfill
\begin{tabular}{lll}
\multicolumn{3}{l}{\textbf{TRENNBARE VERBEN}}\\
& \textbf{ein$|$fügen} & (to insert)\\
& \textbf{hinzu$|$fügen} & (to add)\\
\end{tabular}

\vspace{-5pt}
\end{verbentry}

% führen (to lead / guide / drive)
\begin{verbentry}{
  style = {default},
  toc = {führen (to lead / guide / drive)},
  name = {führen},
  english = {to lead / guide / drive},
  type = {regular},
  auxiliary = {haben}
}
 
    
     \vspace{5pt}

    \hrule
    
    \vspace{5pt}

  \begin{tabular}{rl}
     \multicolumn{2}{l}{\textbf{PRÄSENS} }\\
    ich & \textbf{führe}\\
    du & \textbf{führst}\\
    er/sie/es & \textbf{führt}\\
    wir & \textbf{führen}\\
    ihr & \textbf{führt}\\
    sie/Sie & \textbf{führen}\\
  \end{tabular}
  \hfill
  \begin{tabular}{rl}
    \multicolumn{2}{l}{\textbf{PRÄTERITUM} }\\
    ich & \textbf{führte}\\
    du & \textbf{führtest}\\
    er/sie/es & \textbf{führte}\\
    wir & \textbf{führten}\\
    ihr & \textbf{führtet}\\
    sie/Sie & \textbf{führten}\\
  \end{tabular}
  \hfill
  \begin{tabular}{rl}
     \multicolumn{2}{l}{\textbf{IMPERATIV} }\\
   & \textbf{führe (du)!}\\
   & \textbf{führen wir!}\\
   & \textbf{führt (ihr)!}\\
   & \textbf{führen Sie!}\\
   & \vspace{10pt}
  \end{tabular}

 \vspace{10pt}

 \begin{tabular}{rl}
     \multicolumn{2}{l}{\textbf{PERFEKT}}\\
   & haben + \textbf{geführt}\\
  \end{tabular}
\hfill
\begin{tabular}{rl}
    \multicolumn{2}{l}{\textbf{FUTURE I}}\\
   & werden + \textbf{führen}\\
  \end{tabular}
\hfill
\begin{tabular}{rl}
    \multicolumn{2}{l}{\textbf{PLUSQUAMPERFEKT}}\\
   & hatten + \textbf{geführt}\\
  \end{tabular}

 \vspace{10pt}

\begin{tabular}{lll}
      \multicolumn{3}{l}{\textbf{MIT PRÄPOSITIONEN}}\\
& \textbf{führen zu} + Dat & (to lead to / result in)\\
& \textbf{führen durch} + Akk & (to guide through)\\
& \textbf{führen über} + Akk & (to lead across / over)\\
\end{tabular}
  \hfill
\begin{tabular}{lll}
      \multicolumn{3}{l}{\textbf{TRENNBARE VERBEN}}\\
 & \textbf{an$|$führen} & (to lead to / bring to)\\
 & \textbf{durch$|$führen} & (to guide through)\\
\vspace{10pt}
  \end{tabular}


\vspace{-5pt}
\end{verbentry}

% füllen (to fill)
\begin{verbentry}{
  style = {default},
  toc = {füllen (to fill)},
  name = {füllen},
  english = {to fill},
  type = {regular},
  auxiliary = {haben}
}
 
    
     \vspace{5pt}

    \hrule
    
    \vspace{5pt}

  \begin{tabular}{rl}
     \multicolumn{2}{l}{\textbf{PRÄSENS} }\\
    ich & \textbf{fülle}\\
    du & \textbf{füllst}\\
    er/sie/es & \textbf{füllt}\\
    wir & \textbf{füllen}\\
    ihr & \textbf{füllt}\\
    sie/Sie & \textbf{füllen}\\
  \end{tabular}
  \hfill
     \begin{tabular}{rl}
    \multicolumn{2}{l}{\textbf{PRÄTERITUM} }\\
    ich & \textbf{füllte}\\
    du & \textbf{fülltest}\\
    er/sie/es & \textbf{füllte}\\
    wir & \textbf{füllten}\\
    ihr & \textbf{fülltet}\\
    sie/Sie & \textbf{füllten}\\
  \end{tabular}
  \hfill
  \begin{tabular}{rl}
     \multicolumn{2}{l}{\textbf{IMPERATIV} }\\
   & \textbf{füll(e) (du)!}\\
   & \textbf{füllen wir!}\\
   & \textbf{füllt (ihr)!}\\
   & \textbf{füllen Sie!}\\
   & \vspace{10pt}
  \end{tabular}

    \vspace{10pt}

 \begin{tabular}{rl}
     \multicolumn{2}{l}{\textbf{PERFEKT}}\\
   & haben + \textbf{gefüllt}\\
  \end{tabular}
\hfill
   \begin{tabular}{rl}
    \multicolumn{2}{l}{\textbf{FUTURE I} }\\
   & werde + \textbf{füllen}\\
  \end{tabular}
\hfill
   \begin{tabular}{rl}
      \multicolumn{2}{l}{\textbf{PLUSQUAMPERFEKT}}\\
   & hatten + \textbf{gefüllt}\\
  \end{tabular}

   \vspace{10pt}

   \begin{tabular}{lll}
      \multicolumn{3}{l}{\textbf{MIT PRÄPOSITIONEN}}\\
& \textbf{füllen mit} + Dat & (to fill with)\\
& \textbf{füllen in} + Akk & (to fill into)\\
\end{tabular}
  \hfill
   \begin{tabular}{lll}
      \multicolumn{3}{l}{\textbf{TRENNBARE VERBEN}}\\
 & \textbf{aus$|$füllen} & (to fill out)\\
 & \textbf{auf$|$füllen} & (to refill / top up)\\
\vspace{10pt}
  \end{tabular}

  \vspace{-5pt}
\end{verbentry}

% funktionieren (to function / work)
\begin{verbentry}{
  style = {default},
  toc = {funktionieren (to function / work)},
  name = {funktionieren},
  english = {to function / work},
  type = {regular},
  auxiliary = {haben}
}


\vspace{5pt}
\hrule
\vspace{5pt}

\begin{tabular}{rl}
\multicolumn{2}{l}{\textbf{PRÄSENS}}\\
ich & \textbf{funktioniere}\\
du & \textbf{funktionierst}\\
er/sie/es & \textbf{funktioniert}\\
wir & \textbf{funktionieren}\\
ihr & \textbf{funktioniert}\\
sie/Sie & \textbf{funktionieren}\\
\end{tabular}
\hfill
\begin{tabular}{rl}
\multicolumn{2}{l}{\textbf{PRÄTERITUM}}\\
ich & \textbf{funktionierte}\\
du & \textbf{funktioniertest}\\
er/sie/es & \textbf{funktionierte}\\
wir & \textbf{funktionierten}\\
ihr & \textbf{funktioniertet}\\
sie/Sie & \textbf{funktionierten}\\
\end{tabular}
\hfill
\begin{tabular}{rl}
\multicolumn{2}{l}{\textbf{IMPERATIV}}\\
& \textbf{funktioniere (du)!}\\
& \textbf{funktionieren wir!}\\
& \textbf{funktioniert (ihr)!}\\
& \textbf{funktionieren Sie!}\\
\end{tabular}

\vspace{10pt}

\begin{tabular}{rl}
\multicolumn{2}{l}{\textbf{PERFEKT}}\\
& haben + \textbf{funktioniert}\\
\end{tabular}
\hfill
\begin{tabular}{rl}
\multicolumn{2}{l}{\textbf{FUTURE I}}\\
& werden + \textbf{funktionieren}\\
\end{tabular}
\hfill
\begin{tabular}{rl}
\multicolumn{2}{l}{\textbf{PLUSQUAMPERFEKT}}\\
& hatten + \textbf{funktioniert}\\
\end{tabular}

\vspace{10pt}

\begin{tabular}{lll}
\multicolumn{3}{l}{\textbf{MIT PRÄPOSITIONEN}}\\
& \textbf{funktionieren mit} + Dat & (to function with)\\
\end{tabular}
\hfill
\begin{tabular}{lll}
\multicolumn{3}{l}{\textbf{TRENNBARE VERBEN}}\\
& \textbf{---} & \\
\end{tabular}
\vspace{-5pt}
\end{verbentry}

% füttern (to feed)
\begin{verbentry}{
  style = {default},
  toc = {füttern (to feed)},
  name = {füttern},
  english = {to feed},
  type = {regular},
  auxiliary = {haben}
}
 
    
     \vspace{5pt}

    \hrule
    
    \vspace{5pt}

  \begin{tabular}{rl}
     \multicolumn{2}{l}{\textbf{PRÄSENS} }\\
    ich & \textbf{füttere}\\
    du & \textbf{fütterst}\\
    er/sie/es & \textbf{füttert}\\
    wir & \textbf{füttern}\\
    ihr & \textbf{füttert}\\
    sie/Sie & \textbf{füttern}\\
  \end{tabular}
  \hfill
     \begin{tabular}{rl}
    \multicolumn{2}{l}{\textbf{PRÄTERITUM} }\\
    ich & \textbf{fütterte}\\
    du & \textbf{füttertest}\\
    er/sie/es & \textbf{fütterte}\\
    wir & \textbf{fütterten}\\
    ihr & \textbf{füttertet}\\
    sie/Sie & \textbf{fütterten}\\
  \end{tabular}
  \hfill
  \begin{tabular}{rl}
     \multicolumn{2}{l}{\textbf{IMPERATIV} }\\
   & \textbf{füttere (du)!}\\
   & \textbf{füttern wir!}\\
   & \textbf{füttert (ihr)!}\\
   & \textbf{füttern Sie!}\\
   & \vspace{10pt}
  \end{tabular}

    \vspace{10pt}

 \begin{tabular}{rl}
     \multicolumn{2}{l}{\textbf{PERFEKT}}\\
   & haben + \textbf{gefüttert}\\
  \end{tabular}
\hfill
   \begin{tabular}{rl}
    \multicolumn{2}{l}{\textbf{FUTURE I} }\\
   & werde + \textbf{füttern}\\
  \end{tabular}
\hfill
   \begin{tabular}{rl}
      \multicolumn{2}{l}{\textbf{PLUSQUAMPERFEKT}}\\
   & hatten + \textbf{gefüttert}\\
  \end{tabular}

   \vspace{10pt}

   \begin{tabular}{lll}
      \multicolumn{3}{l}{\textbf{MIT PRÄPOSITIONEN}}\\
& \textbf{füttern mit} + Dat & (to feed with)\\
& \textbf{füttern an} + Akk & (to feed to)\\
\end{tabular}
  \hfill
   \begin{tabular}{lll}
      \multicolumn{3}{l}{\textbf{TRENNBARE VERBEN}}\\
 & \textbf{an$|$füttern} & (to start feeding / bait)\\
 & \textbf{zu$|$füttern} & (to give supplementary food)\\
\vspace{10pt}
  \end{tabular}

  \vspace{-5pt}
\end{verbentry}
